\section{From a Periodic Chain to Energy Bands}
A crystal contains a very large number of repeating microscopic units. Its atomic spacing provides a natural length scale, and the periodicity allows us to describe the one-particle states using a wavevector $k$. Before introducing topology, we construct this description for the simplest possible chain.

\subsection*{1) A tight-binding chain}
Consider one localized orbital $\ket n$ in each unit cell $n$ of a one-dimensional lattice, with lattice spacing $a$. The states are orthonormal, $\braket m n=\delta_{mn}$. An electron can move between neighboring sites by the \emph{hopping amplitude} $t$:
\begin{equation}
 \op H=-t\sum_n\bigl(\ket{n+1}\bra n+\ket n\bra{n+1}\bigr).
\end{equation}
With periodic boundary conditions, $\ket{n+N}=\ket n$. Define the normalized Bloch states
\begin{equation}
\ket k=\frac1{\sqrt N}\sum_{n=0}^{N-1}e^{\ii kna}\ket n,
\qquad k=\frac{2\pi m}{Na},\ m=0,\ldots,N-1.
\end{equation}
Translation by one cell multiplies $\ket k$ by $e^{-\ii ka}$, and the Hamiltonian acts diagonally:
\begin{equation}
 \op H\ket k=-2t\cos(ka)\ket k.
\end{equation}
The dispersion $E(k)=-2t\cos(ka)$ repeats under $k\mapsto k+2\pi/a$. The interval $[-\pi/a,\pi/a)$ is one \emph{Brillouin zone}. A band is the energy of an eigenstate of the periodic Hamiltonian as a function of crystal momentum.

\subsection*{2) More than one orbital per unit cell}
Suppose a cell has two orbitals, $A$ and $B$. The Bloch wavefunction has a two-component coefficient $\ket{u(k)}=(u_A(k),u_B(k))^{\mathsf T}$, and its energies follow from a $2\times2$ Hermitian \emph{Bloch Hamiltonian}:
\begin{equation}
 \op h(k)\ket{u_n(k)}=E_n(k)\ket{u_n(k)},\qquad n=1,2.
\end{equation}
Any $2\times2$ Hermitian matrix can be written
\begin{equation}
 \op h(k)=d_0(k)\mathbf1_2+\vect d(k)\cdot\vect\sigma,\qquad
 E_\pm(k)=d_0(k)\pm|\vect d(k)|.
\end{equation}
The bands are separated by a gap whenever $\vect d(k)\ne0$. At fixed particle filling, a band insulator results if all occupied bands are separated from the empty bands by a finite gap. These are single-particle bands: interactions may require a more general many-body description.

\subsection*{3) The role of geometry}
Besides its energy, a Bloch eigenstate has a direction in the internal space of orbital amplitudes. A smoothly evolving normalized eigenvector $\ket{u_n(k)}$ is defined only up to a $k$-dependent phase:
$\ket{u_n(k)}\mapsto e^{\ii\chi(k)}\ket{u_n(k)}$.
Exactly this freedom appeared in the Berry-phase discussion. As $k$ traverses the Brillouin zone, the eigenvector can accumulate a geometric phase. Certain quantized properties of the eigenvectors remain unchanged when the Hamiltonian is smoothly deformed without closing the gap and without breaking the symmetry that protects the invariant. Such properties are the starting point of topological band theory~\cite{Asboth2016,Zak1989}.